\rhead{SE202: Software Architecture}
\phantomsection 
\section*{Software Architecture (SE202)}
\label{major:SE_2}

\noindent \small{\hyperref[main:scheme]{$\leftarrow$ Back to Scheme}} \vspace{0.5cm}

\noindent \textbf{Credits:} 3 (2-0-2)

\subsection*{Theory}
\noindent\textbf{Unit 1} \\
\textit{Architectural Foundations and Quality Attributes:} Definition and role of software architecture in system development lifecycle, Architectural structures (modules, components, connectors, deployment), Quality attributes (performance, availability, security, maintainability, scalability) and their trade-offs, Architecture Influence Cycles and decision making frameworks, Non-functional requirements elicitation and specification using scenarios and utility trees.

\vspace{0.3cm}

\noindent\textbf{Unit 2} \\
\textit{Architectural Styles and Patterns:} Classification of architectural styles (layered, client-server, pipe-filter, event-driven, service-oriented), Microservices architecture principles and decomposition strategies, Model-View-Controller (MVC) and variations (MVP, MVVM), Architectural patterns (Broker, Blackboard, Peer-to-Peer), RESTful principles and API design best practices for distributed systems.

\vspace{0.3cm}

\noindent\textbf{Unit 3} \\
\textit{Component-Based and Service-Oriented Architectures:} Component principles (cohesion, coupling, interfaces, contracts), CORBA/CCS component model vs. modern container-based approaches, Service-Oriented Architecture (SOA) fundamentals, ESB patterns and anti-patterns, Microservices orchestration vs. choreography, Domain-Driven Design (DDD) bounded contexts and context mapping.

\vspace{0.3cm}

\noindent\textbf{Unit 4} \\
\textit{Architectural Views and Documentation:} 4+1 View Model (Logical, Process, Development, Physical, Scenarios), UML 2.x structure and behavior diagrams for architecture representation, Architecture Description Languages (ADLs), Tool-supported architectural modeling (Enterprise Architect, ArchiMate), Documenting architectural decisions using ADRs (Architecture Decision Records).

\vspace{0.3cm}

\noindent\textbf{Unit 5} \\
\textit{Architecture Evaluation and Evolution:} Architecture Tradeoff Analysis Method (ATAM), Software Architecture Analysis Method (SAAM), Scenario-based evaluation techniques, Risk-driven architecture evaluation, Refactoring legacy architectures, Evolutionary architecture principles (fitness functions, strangler pattern), Cloud-native architecture migration strategies.

\newpage

\subsection*{Practical}
\noindent\textbf{Lab 1: Quality Attribute Scenarios} \\
Focuses on eliciting and documenting non-functional requirements using quality attribute scenarios for a case study system, creating utility trees and prioritizing trade-offs.

\vspace{0.2cm}

\noindent\textbf{Lab 2: Architectural Style Analysis} \\
Focuses on analyzing existing systems to identify applied architectural styles, documenting style characteristics and limitations through reverse engineering exercises.

\vspace{0.2cm}

\noindent\textbf{Lab 3: Microservices Decomposition} \\
Focuses on decomposing monolithic applications into microservices using domain-driven design techniques, identifying bounded contexts and service boundaries.

\vspace{0.2cm}

\noindent\textbf{Lab 4: MVC Implementation Patterns} \\
Focuses on implementing the same business logic using different MVC variations across web frameworks, comparing controller responsibilities and view separation.

\vspace{0.2cm}

\noindent\textbf{Lab 5: REST API Design} \\
Focuses on designing and documenting RESTful APIs following HATEOAS principles, including resource modeling, HTTP methods, status codes, and hypermedia controls.

\vspace{0.2cm}

\noindent\textbf{Lab 6: Component Interface Contracts} \\
Focuses on defining and implementing component interfaces with versioning strategies, demonstrating contract-first design and interface segregation principles.

\vspace{0.2cm}

\noindent\textbf{Lab 7: SOA vs Microservices Simulation} \\
Focuses on implementing service communication patterns (synchronous RPC vs. asynchronous messaging) and measuring performance/scalability differences.

\vspace{0.2cm}

\noindent\textbf{Lab 8: 4+1 View Modeling} \\
Focuses on creating complete 4+1 architectural views for a medium-complexity system using UML tools, ensuring view consistency and scenario coverage.

\vspace{0.2cm}

\noindent\textbf{Lab 9: Architecture Decision Records} \\
Focuses on writing ADRs for architectural decisions in an evolving project, documenting context, decision, consequences, and status changes over time.

\vspace{0.2cm}

\noindent\textbf{Lab 10: ATAM Evaluation Workshop} \\
Focuses on conducting Architecture Tradeoff Analysis Method (ATAM) workshop on peer projects, identifying risks and documenting mitigation strategies.

\vspace{0.2cm}

\newpage
\rhead{Curriculum Schema}
